\section{研究内容与进度}

\begin{frame}{研究内容与进度}{先看一个真实问题：存储系统会丢块}

    \begin{columns}[T,onlytextwidth]
        \begin{column}{0.48\textwidth}
            \begin{block}{为什么不能只靠三副本}
                \footnotesize
                \begin{itemize}
                    \setlength{\itemsep}{3pt}
                    \item 云存储、文件系统和对象存储都要容忍磁盘或节点失效
                    \item 三副本实现简单，但 1\,TB 数据要占约 3\,TB 空间
                    \item 纠删码把数据切成 $k$ 个数据块，再生成 $m$ 个校验块，用更低冗余恢复丢失块
                \end{itemize}
            \end{block}
        \end{column}

        \begin{column}{0.52\textwidth}
            \begin{block}{本课题关心的具体场景}
                \footnotesize
                \begin{itemize}
                    \setlength{\itemsep}{3pt}
                    \item 存储系统中同时可能丢数据块，也可能丢校验块
                    \item 编码和恢复速度越快，系统写入、修复和降级读越不容易被拖慢
                    \item 目标是在 Ceph 这类真实系统里做出高吞吐纠删码实现
                \end{itemize}
            \end{block}

            \vspace{0.3em}
            \begin{block}{一句话概括}
                \centering
                \footnotesize
                少占空间，同时在块丢失时尽快把数据和校验都恢复回来。
            \end{block}
        \end{column}
    \end{columns}

\end{frame}

\begin{frame}{研究内容与进度}{为什么选择移位异或码}

    \setbeamerfont{caption}{size=\scriptsize}
    \begin{columns}[T,onlytextwidth]
        \column{0.42\textwidth}
        \begin{block}{传统 Reed-Solomon 码}
            \scriptsize
            \begin{itemize}
                \setlength{\itemsep}{1pt}
                \item 可靠、成熟，Ceph 等系统已广泛使用
                \item 编解码依赖有限域矩阵运算
                \item 高性能实现优化复杂
            \end{itemize}
        \end{block}

        \begin{block}{Two-tone 移位异或码}
            \scriptsize
            \begin{itemize}
                \setlength{\itemsep}{1pt}
                \item 校验块由数据块移位后异或得到
                \item 移位可用地址偏移实现
                \item 理论上很快，但缺少工程级高性能实现
            \end{itemize}
        \end{block}

        \column{0.58\textwidth}
        \centering
        \scriptsize\textbf{例：4 个数据块经过不同偏移后异或生成 3 个校验块}

        \vspace{0.2em}
        \includegraphics[width=0.96\textwidth,height=0.49\textheight,keepaspectratio]{Example of Two-tone Encoding.pdf}
    \end{columns}

\end{frame}

\begin{frame}{研究内容与进度}{中期已经做出了什么}

    \begin{columns}[T,onlytextwidth]
        \column{0.50\textwidth}
        \begin{block}{已经完成的系统}
            \footnotesize
            \begin{itemize}
                \setlength{\itemsep}{3pt}
                \item 完成 Two-tone 编码、解码和修复流程分析
                \item 实现约 2000 行 C++ FastTT 核心库
                \item 把 FastTT 做成 Ceph v18.2.7 纠删码插件
                \item 通过 Ceph 接口测试和集群端到端读写验证
            \end{itemize}
        \end{block}

        \column{0.50\textwidth}
        \begin{block}{已经得到的结果}
            \footnotesize
            \begin{itemize}
                \setlength{\itemsep}{3pt}
                \item 相较 ISA-L：编码平均提升 \textbf{59.1\%}
                \item 相较 ISA-L：解码平均提升 \textbf{50.9\%}
                \item 单块丢失恢复最高提升 \textbf{182.1\%}
                \item 与学术基线 Cerasure 对比也取得整体优势
                \item 相关成果已整理为论文并投稿 IPDPS 2026
            \end{itemize}
        \end{block}
    \end{columns}

\end{frame}
